\documentclass[10pt,a4paper]{amsart}
\usepackage[margin=19mm]{geometry}
\usepackage{amsmath,amssymb,mathtools}
\usepackage[scr=rsfs,cal=euler]{mathalfa}
\usepackage{xfrac}
\usepackage[unicode,hidelinks,bookmarks=false]{hyperref}
\newcommand{\ie}{i.e.\ }
\newcommand{\cf}{cf.\ }
\newcommand{\resp}{respectively\ }
\newcommand{\Subsection}{\subsection}
\providecommand{\nobreakdash}{\nobreak\hbox{-}}
\setlength{\emergencystretch}{2em}
\multlinegap0pt
\usepackage{amssymb}
\usepackage{tikz}
\newtheorem{lemma}{Lemma}
\newcommand{\ver}{\mathrm{ver}}
\newcommand{\vol}{\mathrm{vol}}
\title{Spectra of normalized volumes \\ of right-angled hyperbolic polyhedra}
\author{A. Egorov, A. Vesnin}
\date{Source: arXiv:2605.18558v2 (CC BY 4.0)}
\begin{document}
\maketitle

\noindent\emph{Source and licence.} Spectra of normalized volumes \\ of right-angled hyperbolic polyhedra. Version arXiv:2605.18558v2 (CC BY 4.0), distributed by arXiv under the Creative Commons Attribution 4.0 International licence, http://creativecommons.org/licenses/by/4.0/, as recorded in the arXiv licence metadata for this exact version and shown on the abstract page https://arxiv.org/abs/2605.18558v2. Full licence notice: \texttt{licenses/arxiv-2605.18558-CC-BY-4.0.txt}.\par\smallskip

\noindent\textit{Editorial scope.} Counted unit: the article's lemma lemma4.1 (source0.tex lines 1303--1313). The article's complete original proof of it is delivered with it (source0.tex lines 1314--1317, one \texttt{proof} environment of 4 lines). No mathematics is written, repaired or re-derived: the statement and the proof are the article's own lines, in the article's own notation, and no retained line is substituted or re-flowed.  No further source-local context is needed: every cross-reference of every retained line resolves inside the two delivered spans.  Nothing was omitted from any retained span, so the package carries no omission record.  No retained line contains a citation, so the wrapper carries no bibliography and no bibliography entry.  No retained line contains a figure environment, an image-inclusion command or drawing code, so no image is delivered and the package declares no asset dependency.  Editorial formatting only, in addition: an AMS \texttt{amsart} preamble that restates the article's own declarations and macros used by the retained lines, namely the environments \texttt{lemma} on separate counters, together with the standard packages those lines need.  The retained environments print in this document as Lemma 1 in order of appearance, whereas the article prints the same items with the numbering of its own full document.  The complete native source of the article as distributed in the arXiv e-print for this exact version is bundled unchanged as \texttt{sources/200863/source0.tex}.
	\begin{lemma} \label{lemma4.1}
		Under the operation of connected sum of ideal right-angled polyhedra along triangular faces, volume behaves additively: 
		$$
		\vol(\mathcal P_1 \# \mathcal P_2) = \vol(\mathcal P_1) + \vol(\mathcal P_2),
		$$
		and the number of vertices changes as follows:
		$$
		\ver(\mathcal P_1 \# \mathcal P_2) = \ver(\mathcal P_1) + \ver(\mathcal P_2) - 3.
		$$
		Moreover, the polyhedron $\mathcal P_1 \# \mathcal P_2$ is ideal right-angled.
	\end{lemma}

	\begin{proof}
			When right-angled polyhedra are glued, the faces adjacent to $T_1$ and $T_2$ pairwise form flat dihedral angles of size $\pi$ and merge into complete geodesic hyperfaces. The right angles with the remaining faces are preserved; hence $\mathcal P_1 \# \mathcal P_2$ is again a right-angled polyhedron.
			Since the two-dimensional gluing face has zero three-dimensional volume, the equality $\vol(\mathcal P_1 \# \mathcal P_2) = \vol(\mathcal P_1) + \vol(\mathcal P_2)$ holds. When the gluing faces are identified, the three ideal vertices of the triangle $T_1$ are identified pairwise with the three ideal vertices of the triangle $T_2$. Since these vertices lie on the absolute, after gluing they remain ideal vertices. Thus, the $6$ original vertices corresponding to $T_1$ and $T_2$ merge into $3$ vertices, which gives the formula $\ver(\mathcal P_1 \# \mathcal P_2) = \ver(\mathcal P_1) + \ver(\mathcal P_2) - 3$. The lemma is proved.
	\end{proof}

\end{document}
